% !TeX spellcheck = en_us
\documentclass[a4paper,12pt,fleqn]{article}

\usepackage[left=1.50cm, right=1.50cm, top=3cm, bottom=2cm, 
headheight= 2.5cm, headsep= 4pt, footskip=25pt]{geometry}
\usepackage{fancyhdr}
\usepackage{lmodern}
\usepackage[table]{xcolor}
\usepackage{amssymb,amstext,amsmath}
\usepackage{booktabs}
\usepackage{bm}
\usepackage{setspace}
\usepackage{graphicx}
\usepackage[colorlinks]{hyperref}
\setlength{\parindent}{0mm}
\setlength{\parskip}{2mm}
\usepackage{enumitem}
\usepackage{titlesec}
\usepackage{gensymb}
\usepackage{pgfgantt}
\usepackage{lscape}
\usepackage[absolute]{textpos}
\usepackage{comment}
\usepackage{cite}
\usepackage{tikz}
\usepackage{pgfplots}
\usetikzlibrary{shapes.geometric}
\usepackage{array}
\newcolumntype{P}[1]{>{\centering\arraybackslash}p{#1}}
\usepackage{longtable}
\usepackage[font={small}]{caption}
\usepackage[most]{tcolorbox}
\usepackage{tikz}
\usepackage{stmaryrd}
\usepackage{wrapfig}
%\usepackage[none]{hyphenat}
\usepackage{ctable}
\usepackage{ragged2e}
\usepackage[official]{eurosym}
\newcommand*\circled[1]{\tikz[baseline=(char.base)]{
		\node[shape=circle,draw,inner sep=1pt] (char) {#1};}}
\usepackage[symbol]{footmisc}
\renewcommand{\thefootnote}{\fnsymbol{footnote}}
\renewcommand{\thempfootnote}{\fnsymbol{mpfootnote}}
\renewcommand*\footnoterule{}

\usepackage{xcolor}
\definecolor{docColor}{cmyk}{1.00	, 1.00  , 0   , 0.40}  % 0.90\% of black
\color{docColor}

\newtcolorbox{myhlbox}[2][]{colback=white, colframe=docColor, width=1.0\linewidth,
	boxrule=1.0pt, left=0.0pt, right=0.0pt, top=4pt, bottom=2.0pt,
	colbacktitle=docColor!20, coltitle=black, enhanced, attach boxed title to top left={xshift=3mm, yshift=-1mm},
	title=#2, #1}

\hypersetup{
	colorlinks=true,
	citecolor=blue,
	filecolor=magenta,
	linkcolor=blue,
	urlcolor=blue,
	pdftitle={SeminarAnnouncements},
	%bookmarks=true,,
}

\graphicspath{{./figs/}}

\definecolor{dgreen}{rgb}{0.0, 0.5, 0.0}

\allowdisplaybreaks

\newcommand*\mystrut[1]{\vrule width0pt height0pt depth#1\relax}
\renewcommand{\thesection}{\arabic{section}.}
\renewcommand{\thesubsection}{\arabic{section}.\arabic{subsection}}
\renewcommand{\thesubsubsection}{\arabic{section}.\arabic{subsection}.\arabic{subsubsection}}
\titleformat{\section}{\sffamily\large\bfseries}{\thesection}{1.0em}{}
\titleformat{\subsection}{\sffamily\large\bfseries}{\thesubsection}{1.0em}{}
\titleformat{\subsubsection}{\sffamily\normalsize\bfseries}{\thesubsubsection}{1.0em}{}
\titlespacing\section{0pt}{0pt plus 2pt minus 2pt}{0pt plus 2pt minus 2pt}
\renewcommand{\baselinestretch}{1.00}

\newcommand\CircledImage[1]{%
	\tikz\path[fill overzoom image={#1}, draw=docColor, line width=0mm]circle[radius=0.5\linewidth];%
}

\input{format/defs}

\fancypagestyle{titlepage}{
	\fancyhf{}
	\insertLXheader
}

\fancypagestyle{plain}{%
	\fancyhf{}
	\renewcommand{\headrulewidth}{0pt}
	\renewcommand{\footrulewidth}{0pt}
	\insertnormalpage
}

\pagestyle{plain}

\begin{document}
	\sffamily
\thispagestyle{titlepage}
\vspace*{-3em}
\begin{center}
	\huge \textbf{LMS Seminar}
\end{center}
\begin{center}
	\Large 
	\textbf{Factors controlling VHCF life of Ni-based single crystal superalloys} 
\end{center}
\begin{center}
	\Large
	Jonathan Cormier \\
	{\large Institut PPRIME, CNRS \\ 
	Universit{\'e} de Poitiers, ISAE-ENSMA} \\
	\vspace*{1em}
	{\Large
	\textbf{Date and Time}:~October 02, 2025 (2 -- 3 pm)
	
	\textbf{Venue}:~Amphi 104 (Pole Meca) }
\end{center}
\vspace*{-1em}
\begin{myhlbox}{{\large Abstract}}
	Ni-based single crystal superalloys (SX) are widely used in gas turbine engines for the manufacturing of high pressure turbine blades due to their exceptional mechanical properties at high temperature. Service operations of blades may lead to fatigue controlled failure mechanisms due to the vibrations introduced by the gas flow in addition to the centrifugal forces. These failures are difficult to forecast, as up to 90\% of the fatigue life is spent in the crack initiation phase. Typical frequencies of vibrations of airfoils are in between 1 and 10 kHz, requiring specific experimental facilities to achieve the very high cycle fatigue (VHCF) domain at high temperature.
	
	In this presentation, a critical analysis of the VHCF life sensitivity at 1,000 deg. C/20 kHz, R = -1 and R $>$ 0.3 to the processing parameters (dendritic chemical homogeneity, casting pore size, introduction of a prior plastic deformation and $\gamma/\gamma$’ microstructure degradation - i.e. $\gamma$’ rafting) will be performed. For this, 10 different Ni-based SX alloys have been investigated, with different chemical compositions and/or processing parameters. A special attention will be paid in this presentation on the fine scale crack initiation mechanisms, involving deformation mechanisms far from being the usually observed ones at slow strain rates at this temperature (strain rates $<$ 10-2 s-1). By varying the casting process or the oxidation resistance among the different alloys studied, a map of crack initiation mechanisms in VHCF will finally be proposed for Ni-based SX alloys. Finally, some insights on the role of a coating will also be presented.
	
	
\end{myhlbox}

\vspace*{1em}

\begin{myhlbox}{{\large About the speaker}}
	\begin{wrapfigure}{r}{0.30\textwidth}
		\vspace*{0em}
		\CircledImage{images/cormier}
	\end{wrapfigure}
	Jonathan Cormier is a professor in the Physics and Mechanics of Materials department at ISAE-ENSMA. He received his Engineering diploma and PhD degree at the same institute in 2003 and 2006, respectively. Prof. Cormier is an expert on a wide range of problems related to the high-temperature behavior and durability of materials with specific interests in superalloys, thermal barrier coatings and ceramic matrix composites among others. He holds the EMERAUDE chair, a large research program between  Safran and ISAE-ENSMA focusing on processing-microstructure-properties relationships of Ni-based superalloys. He is the recipient of multiple awards e.g., the Jean-Rist medal for metallurgy, FEMS lecturer award for excellence in materials science and engineering, to name a couple. He also serves as editor-in-chief of the journal Metallurgical and Materials Transactions A.
	
\end{myhlbox}
\end{document}